\section{\texorpdfstring{Levels in $\mathcal C$ and positive exception mass on $\mathcal G$}{Levels in C and positive exception mass on G}}\label{sec:common}

We use the histories, history lists and three edge-pair scores constructed in
\Cref{sec:budgets}. The first objective is to show that levels in $\mathcal C$ contribute
only \(o(S)\) bit variation at a center. The second is to show that the point
exceptions on levels in $\mathcal G$ retain a positive fraction of \(S\). Both
conclusions come from comparing the upper and lower edge-pair bounds, but the
comparison first requires a precise description of the states permitted by a
history list.

Throughout this section, \(S\to\infty\), the parameters satisfy
\Cref{bud:parameters}, and \(H\) is the fixed history map from
\Cref{bud:prediction-error}. At a level, \(\nu_h\) denotes the distribution of
the true state \((A,B_1)\) of a point sampled from \(\mu\) conditional on
\(H=h\). Dependence on the level is suppressed. Recall from
\eqref{pred:strict-cross} that the states agreeing with \(y=(a,b)\) in exactly one coordinate form
\[
 R_y=\set{(u,v):(u=a,\ v\ne b)\ \text{or}\ (u\ne a,\ v=b)}.
\]
A row means a set with fixed first coordinate; a column has fixed second
coordinate.

\begin{proposition}[Bit disagreement on $\mathcal C$]\label{com:variation}
Under the standing parameter and history choices, the bit-disagreement integrals for experiments $Q$ and $P\theta$ from \eqref{bud:variations} satisfy
\begin{equation}\label{com:variation-eq}
 C_Q+C_{P\theta}=o(S).
\end{equation}
Moreover their exception means on $\mathcal G$ satisfy
\begin{equation}\label{com:center-rare}
 \E_\kappa H_q^*\le\E_\mu H_p^*+o(S).
\end{equation}
\end{proposition}

\begin{proposition}[Positive exception mass on $\mathcal G$]\label{com:rare-positive}
For the same incidence subgraphs, parameters and fixed history, the point exceptions on $\mathcal G$ satisfy
\begin{equation}\label{com:rare-positive-eq}
 \liminf_{t\to\infty}\frac{\E_\mu H_p^*}{S}>0.
\end{equation}
\end{proposition}

The proof of \Cref{com:variation} is in
\Cref{sec:common-variation-proof}. On levels in $\mathcal C$ the state
classification gives a lower bound with coefficient \(\varpi/4\)
for the sum of the two variations. The terms from outside $\mathcal C$ cancel between
the two orientations by \Cref{bud:rare-lower}, whereas
\Cref{bud:upper} bounds the same sum with coefficient
\(10\delta^2\varpi<\varpi/4\). The remaining proposition is
proved in \Cref{sec:rare-positive-proof}: if the point-exception mean
were \(o(S)\), a second comparison would also give \(C_P=o(S)\),
and the matching-label estimate would force \(S=o(S)\).

\subsection{From sampled histories to the actual point}

For a center \(q\), define the mixture of predictions
\[
 \widetilde\nu_q=\E_{p\mid q}\nu_{H(p)}.
\]
For its history list \(\theta=(h_1,\ldots,h_k)\), define the finite set of candidate
center states
\begin{equation}\label{com:candidates}
 \mathcal Y_\theta
 =\set{y\text{ a center state}:\nu_{h_j}(R_y)\ge1-k^{-2}
                      \text{ for every }1\le j\le k}.
\end{equation}
The actual state \(\mathbf y_q\) is usually a candidate: the union bound and Markov's
inequality give, after integration on \(\mathcal C\),
\begin{equation}\label{com:actual-candidate}
 \int_{\mathcal C}\Prob(\mathbf y_q\notin\mathcal Y_\theta)
 \le k^3\int_{\mathcal C}
       \E_E[1-\nu_{H(p)}(R_{\mathbf y_q})]=o(S).
\end{equation}
We also need \(R_y\) for every candidate \(y\) to have large mass under
\(\widetilde\nu_q\). The following elementary sampling bound supplies this
uniformity.
For its first assertion, we specialize the \(\varepsilon\)-net
double-sampling argument to the sets \(R_y\)
\cite[Lemmas~3.4--3.5]{HausslerWelzl1987}.
The trace bound for these sets, the weighted-law formulation, and the subsequent
transfer of concentration of the conditional law are established here.

\Needspace{5\baselineskip}
\begin{lemma}[Sampling the sets \(R_y\)]\label{com:ghost}
Let \(\nu\) be a probability measure on a finite product set, let
\(k\ge1\) be an integer, and let \(X_1,\ldots,X_k\) be independent
samples from \(\nu\). For
\(0<\alpha<1\) and \(k\alpha\ge8\),
\[
 \Prob\bigl(\exists y:\nu(R_y)<1-\alpha,
                    \ X_1,\ldots,X_k\in R_y\bigr)
 \le 2(2k+1)^2\,2^{-k\alpha/2}.
\]
Consequently, when \(k\) is sufficiently large, the history lists in
\eqref{com:candidates} satisfy
\begin{equation}\label{com:mixture}
 \Prob\bigl(\exists y\in\mathcal Y_\theta:
              \widetilde\nu_q(R_y)<1-k^{-1/4}\mid q\bigr)
 \le 4(2k+1)^2\,2^{-k^{3/4}/2}
\end{equation}
at each level, uniformly in \(q\).
\end{lemma}

\begin{proof}
Adjoin independent samples \(X_{k+1},\ldots,X_{2k}\). If the first \(k\)
samples admit a set \(R_y\) as in the event, select one by a fixed ordering. For this
selected set \(R_y\), the number of misses among the new samples has mean
\(k\beta\), where \(\beta=\nu(R_y^c)>\alpha\), and variance at most
\(k\beta\). Thus the probability of fewer than \(k\alpha/2\) misses is at
most \(4/(k\beta)\le1/2\). It suffices to bound twice the probability that
some \(R_y\) contains the first half and misses at least \(k\alpha/2\) of the
second half.

The joint sample mass is \(\prod_{i=1}^{2k}\nu(X_i)\), which is unchanged
by every permutation of the indices, even when \(\nu\) is not uniform.
We may therefore condition on the pooled \(2k\) samples and split their
indices uniformly into two halves. A set \(R_y\) induces at most \((2k+1)^2\) different subsets of these
indices: its first coordinate can agree with one of at most \(2k\) observed
coordinates, or none, and the same holds for its second coordinate. If a
pattern has \(m\) misses, the probability that they all fall in the second
half is zero for \(m>k\), and otherwise is
\[
 \frac{k(k-1)\cdots(k-m+1)}{(2k)(2k-1)\cdots(2k-m+1)}
 \le 2^{-m}.
\]
The union bound proves the first assertion. Repeated sampled states cause no
problem, since the argument concerns the indexed samples.

For the second assertion, after sampling the histories independently given
\(q\), independently sample a state from each corresponding \(\nu_{h_j}\).
Before the histories are revealed these states are independent with law
\(\widetilde\nu_q\). If a candidate violating the required conclusion exists,
select one as a function of the history list. Conditional on that history list, all
the synthetic states lie in its set \(R_y\) with probability at least
\(1-k\cdot k^{-2}=1-1/k\). Apply the first assertion with
\(\alpha=k^{-1/4}\) and divide by \(1-1/k\).
\end{proof}

Since \(W_c\le k^{1/16}S\), the integral of the failure probability in
\eqref{com:mixture} is \(o(S)\). A candidate therefore constrains the mixture
of predicted states. Concentration of this mixture does not by itself
control the actual point: a set chosen from \((q,\theta)\) may use
information about the point that its history does not contain.
For rows, columns and sets of at most two states, the next lemma
transfers this concentration to the actual point using exceptional sets
determined only by the history and level.

\begin{lemma}[Transfer for restricted sets]\label{com:transfer}
Fix a level and the joint edge--history list law
\[
 \Prob(p,q,\theta)=\Prob_E(p,q)\Prob(\theta\mid q).
\]
Let \(0<\eta<1/2\). As a function of \((q,\theta)\), select a set
\(U(q,\theta)\) that is a row, a column, or a set of at most two states.
On an event \(G\) determined by \((q,\theta)\), suppose
\[
 \widetilde\nu_q(U(q,\theta))\ge1-3\eta^2.
\]
Then the actual point state satisfies
\[
 \Prob\bigl(G,\ \mathbf y_p\notin U(q,\theta)\bigr)\le8\eta.
\]
The set type and its coordinates may depend on both \(q\) and \(\theta\).
\end{lemma}

\begin{proof}
Conditional on \((q,\theta)\), the history of the actual point has its
ordinary conditional law given \(q\), because the history list is independent of
the point given \(q\). Hence on \(G\),
\[
 \E[\nu_{H(p)}(U(q,\theta)^c)\mid q,\theta]\le3\eta^2.
\]
The event that this set has \(\nu_{H(p)}\)-mass less than \(1-\eta\)
therefore has probability at most \(3\eta\).

For each history \(h\), define three exceptional sets using only \(\nu_h\).
If a row has mass at least \(1-\eta\), it is unique, and the first exceptional
set is its complement. If no such row exists, the first exceptional set is
empty. Define the second set in the same way for columns. Their masses are
at most \(\eta\) each.

If a set of at most two states has mass at least \(1-\eta\), put
\[
 M_h=\set{x:\nu_h(x)>\eta}
\]
and take \(M_h^c\) as the third exceptional set. Every set \(U\) of at most
two states and mass at least \(1-\eta\) contains \(M_h\): omitting an atom
of mass greater than \(\eta\) would violate that mass bound. Choosing one
such \(U\) also gives
\[
 \nu_h(M_h^c)\le\nu_h(U^c)+\nu_h(U\setminus M_h)
                 \le\eta+2\eta=3\eta.
\]
If no such \(U\) exists, the third exceptional set is empty. The union of
the three exceptional sets thus has mass at most \(5\eta\).

The actual point has marginal law \(\mu\), under which its state conditional
on \(H=h\) has distribution exactly \(\nu_h\). Consequently the probability
of this union of exceptional sets is at most \(5\eta\) under the entire
edge--history list law. Outside it and the earlier nonconcentration event, the
actual state belongs to whichever permitted set was selected. Adding
the two error bounds proves the result.
\end{proof}

In the application, at most three candidates will be chosen by a fixed
ordering at each level and history list. Outside the failure event bounded in
\eqref{com:mixture}, the intersection of their sets \(R_y\) has mixture mass at
least \(1-3k^{-1/4}\). The set will be this intersection or a row or
column containing it. Its choice, the transfer event and its cost are
specified with the bad edges in \Cref{sec:common-variation-proof}.

\subsection{\texorpdfstring{Classification and suppression of bit disagreement on $\mathcal C$}{Classification and suppression of bit disagreement on C}}\label{sec:common-variation-proof}

We now apply this transfer to the intersections of the sets \(R_y\). Two
distinct center states in one row, say \((a,b)\) and \((a,b')\), give the intersection
\begin{equation}\label{com:row-intersection}
 R_{(a,b)}\cap R_{(a,b')}
       =\set{(a,v):v\notin\set{b,b'}}.
\end{equation}
A third center state outside that row cuts this intersection down to at most
one point. The corresponding assertions hold for columns. If two states
\((a,b)\) and \((a',b')\) differ in both coordinates, the intersection of their sets is
\begin{equation}\label{com:corners}
 R_{(a,b)}\cap R_{(a',b')}
       =\set{(a,b'),(a',b)}.
\end{equation}
Three states whose first coordinates are all distinct and whose second
coordinates are all distinct have empty intersection of their sets \(R_y\). Indeed each of the
two points in \eqref{com:corners} fails to share a coordinate with the third
state. These observations give an exhaustive classification: a candidate
set is empty, is a singleton, contains two distinct states sharing a
coordinate, or consists of states with pairwise distinct coordinates.
In the shared-coordinate case, either all candidates lie in that row or
column, or a third lies outside it. In the pairwise-distinct case there
are either two candidates or at least three.

\Cref{com:fig-crosses} shows the two-corner case: either endpoint state
together with the bit determines the state at the other endpoint.
\begin{figure}[tb]
\centering
\begin{tikzpicture}[x=1.45cm,y=1.15cm]
  \draw[gray!45] (-.55,0)--(2.55,0) (-.55,2)--(2.55,2);
  \draw[gray!45] (0,-.5)--(0,2.5) (2,-.5)--(2,2.5);
  \node[draw,circle,fill=white,inner sep=2.2pt] (y1) at (0,0) {};
  \node[draw,circle,fill=white,inner sep=2.2pt] (y2) at (2,2) {};
  \node[circle,fill=blue!65!black,inner sep=2.5pt] (x1) at (2,0) {};
  \node[circle,fill=blue!65!black,inner sep=2.5pt] (x2) at (0,2) {};
  \node[below left=3pt of y1] {\(y_1=(a,b)\)};
  \node[above right=3pt of y2] {\(y_2=(a',b')\)};
  \node[below right=3pt of x1] {\((a,b')\)};
  \node[above left=3pt of x2] {\((a',b)\)};
  \node[rotate=90,fill=white,inner sep=2pt] at (0,1) {\(Y=0\)};
  \node[rotate=90,fill=white,inner sep=2pt] at (2,1) {\(Y=0\)};
  \node[fill=white,inner sep=2pt] at (1,0) {\(Y=1\)};
  \node[fill=white,inner sep=2pt] at (1,2) {\(Y=1\)};
\end{tikzpicture}
\caption{Two candidate center states differing in both coordinates admit
exactly two point states in the intersection of their sets \(R_y\). The filled points are
those opposite corners. Within this four-state configuration, the bit of a
consistent edge and the state at either endpoint determine the state at
the other endpoint. Rows have fixed first label and columns have fixed
second label. The diagram displays coordinate
equalities of states.}
\label{com:fig-crosses}
\end{figure}


For each history list with at least two candidates, choose witnesses by fixed
orderings: two states sharing a coordinate, with a third outside that row
or column when one exists; otherwise two candidates if there are exactly
two, and three if there are more. If all candidates lie in the row or
column of the shared-coordinate pair, let \(U_\theta\) be that row or
column. In the other cases let \(U_\theta\) be the intersection of the sets \(R_y\) for the witnesses, which has at most two states. Set \(U_\theta=\varnothing\)
for empty or singleton candidate sets.

Let \(G\) be the event, determined by \((q,\theta)\) at this level, that
\(\abs{\mathcal Y_\theta}\ge2\) and
\(\widetilde\nu_q(R_y)\ge1-k^{-1/4}\) for every
\(y\in\mathcal Y_\theta\). On \(G\), the at most three witness sets \(R_y\)
give \(\widetilde\nu_q(U_\theta)\ge1-3k^{-1/4}\). Thus
\Cref{com:transfer}, with \(\eta=k^{-1/8}\), gives
\[
 \Prob(G,\ \mathbf y_p\notin U_\theta)\le8k^{-1/8}.
\]
Declare \((p,q,\theta)\) bad if its edge is inconsistent, if
\(\mathbf y_q\notin\mathcal Y_\theta\), if the failure event in
\eqref{com:mixture} occurs, or if \(G\) holds and \(\mathbf y_p\notin U_\theta\).
Empty candidate sets force the second condition. For large \(k\), empty
witness intersections force the third, since their mixture mass cannot
be at least \(1-3k^{-1/4}>0\); hence they have no good edges either.
The transfer event contributes at most
\(8k^{-1/8}W_c\le8k^{-1/16}S=o(S)\). Together with the preceding
estimates and the integrated inconsistency bound, this gives
\begin{equation}\label{com:bad-mass}
 \int_{\mathcal C}\Prob(p,q,\theta\text{ is bad})=o(S).
\end{equation}

The comparisons below must use the full conditional bit probability
\(\rho\) in the definition of the scores. We never remove bad edges and
renormalize that law. In particular the discontinuity of \(c(\rho)\) at
\(1/2\pm\delta\) introduces no error.

Here is the numerical estimate used in the two-corner case. For every
\(0\le\rho\le1\), with \(\delta=0.01\), \(\varpi=\delta^2/2\), and
\(c\) as in \Cref{sec:budgets},
\begin{align}
 c(\rho)\bigl(\rho^2+(1-\rho)^2\bigr)-2\rho(1-\rho)
 &= (2\rho-1)^2+
       (c(\rho)-1)\bigl(\rho^2+(1-\rho)^2\bigr)\notag\\
 &\ge\varpi/2.\label{com:score-polynomial}
\end{align}
For \(\abs{\rho-1/2}\le\delta\), discard the first square and use
\(\rho^2+(1-\rho)^2\ge1/2\). Outside that interval the expression is at
least \(4\delta^2-\varpi>\varpi/2\).

\Needspace{5\baselineskip}


\begin{proof}[Proof of \Cref{com:variation}]
Fix a level in $\mathcal C$. For each history list of positive probability, let
\(w_\theta=\Prob(\theta)\) and
\(b_\theta=\Prob(\text{bad edge}\mid\theta)\), with both probabilities taken
under the joint law in \Cref{com:transfer}. In the experiment $Q$,
condition on \(\theta\) and then sample \(q\) by its actual conditional law.
The two points are still independent with their ordinary law given \(q\).
In the experiment $P\theta$, sample \(p\) conditional on
\(\theta\), then two independent centers given \((p,\theta)\). Each slot
in either experiment has the same conditional single-edge law
\(\Prob(p,q\mid\theta)\). Therefore a pair with at least one bad slot has
probability at most \(2b_\theta\) in either experiment.

Write \(V_Q(\theta)\) and \(V_{P\theta}(\theta)\) for the corresponding
averages of \(2\rho(1-\rho)\) over bases in this slice, and write
\(d_Q(\theta),d_{P\theta}(\theta)\) for its two expected scores on $\mathcal C$. Each \(V\) is at most \(1/2\), and either score is at least the
negative of its variation. We prove
\begin{equation}\label{com:slice-score}
 d_Q(\theta)+d_{P\theta}(\theta)
 \ge\frac\varpi4\bigl(V_Q(\theta)+V_{P\theta}(\theta)\bigr)
       -C b_\theta
\end{equation}
with an absolute constant \(C\). All coefficients in a score on $\mathcal C$ are
bounded by \(1+\varpi\), so charging a pair with a bad slot always costs
at most an absolute multiple of its probability.

\emph{An empty candidate set.}
Every edge is bad, so \(b_\theta=1\), and the boundedness of the scores proves
\eqref{com:slice-score} on increasing \(C\).

\emph{One candidate state.}
At a fixed base $(p,\theta)$, two good edges have the same center state.
Consistency forces their bits to agree, and the nonshared coordinate of
their center states matches. Thus this edge-pair score is at least
\(1-\varpi-O(b_\theta)\) after averaging over its bases. The score for an edge pair at a center is at least \(-1/2\). Their sum is at least
\(1/2-\varpi-O(b_\theta)\), which implies \eqref{com:slice-score} because
the two variations sum to at most one.

\emph{Two candidates sharing a coordinate, and all candidates in that row
or column.}
The transfer using \eqref{com:row-intersection} places every good point in
the same row or column. If it is a row, every good edge has matching
\(A\)-coordinates and hence bit \(1\), by consistency; in the column case
every good edge has bit \(0\). A bit difference in either pair experiment
therefore requires a bad slot. This gives
\(V_Q(\theta)+V_{P\theta}(\theta)\le4b_\theta\). The lower score bound
\(d_i\ge-V_i\) now proves \eqref{com:slice-score}.

\emph{Two candidates sharing a coordinate, and a third outside that row
or column.}
The three witness sets \(R_y\) intersect in at most one state. If the
intersection is empty, all edges are bad. Otherwise all good point states
equal that one state. At a center base, two good edges consequently have
equal bits and the required nonshared-coordinate match. Its score is at
least \(1-\varpi-O(b_\theta)\), and the other score is at least \(-1/2\).
The same estimate as in the singleton case applies.

\emph{No coordinate shared by distinct candidates.}
If there are at least three candidates, their sets \(R_y\) have empty intersection and
all edges are bad. Otherwise there are exactly two candidates. The good
point states lie in their two opposite corners from \eqref{com:corners}.
At either type of base, a good edge's bit determines the state at the other
endpoint. Thus two good edges with equal bits have the required collision.
If \(\beta_x\) is the conditional probability of a bad slot at the base,
then under its full pair law
\[
 \Prob(Y_1=Y_2,N\mid\text{base})
 \ge \rho^2+(1-\rho)^2-2\beta_x.
\]
The bit-disagreement probability is still exactly \(2\rho(1-\rho)\).
Thus the expected score is at least the left side of
\eqref{com:score-polynomial} minus \(2(1+\varpi)\beta_x\).
Averaging \(\beta_x\) over either type of base gives \(b_\theta\). Hence
each score is at least \(\varpi/2-O(b_\theta)\), which proves
\eqref{com:slice-score}.

This completes the classification. Multiply \eqref{com:slice-score} by its
actual slice weight \(w_\theta\), sum, and integrate over levels in $\mathcal C$.
The error is
\(C\int_{\mathcal C}\sum_\theta w_\theta b_\theta=o(S)\) by
\eqref{com:bad-mass}. This argument allows individual \(b_\theta\) to be
large; no inverse slice weight or number of history lists enters. The contributions on $\mathcal C$ to \(\mathcal D_Q+\mathcal D_{P\theta}\) are therefore at least
\[
 \frac\varpi4(C_Q+C_{P\theta})-o(S).
\]
The lower bounds in \Cref{bud:rare-lower} outside $\mathcal C$ cancel their
main terms, because the two base vertices are on opposite sides. Hence
\[
 \frac\varpi4(C_Q+C_{P\theta})-o(S)
 \le \mathcal D_Q+\mathcal D_{P\theta}
 \le 10\delta^2\varpi(C_Q+C_{P\theta})+o(S),
\]
where the last inequality is \Cref{bud:upper}. Since
\(10\delta^2<1/4\), this proves \eqref{com:variation-eq}.

Finally the contribution on $\mathcal C$ to \(\mathcal D_Q\) is at least \(-C_Q\).
\Cref{bud:rare-lower} and \Cref{bud:upper} therefore give
\[
 2\E_\kappa H_q^*-2\E_\mu H_p^*-C_Q-o(S)
 \le \mathcal D_Q\le10\delta^2\varpi C_Q+o(S).
\]
Use \(C_Q=o(S)\) to obtain \eqref{com:center-rare}.
\end{proof}

\subsection{\texorpdfstring{A second score forces positive exception mass on $\mathcal G$}{A second score forces positive exception mass on G}}\label{sec:rare-positive-proof}

We have suppressed bit disagreement on $\mathcal C$ at centers. To show that point exceptions
remain on levels with $f<\epsilon$, suppose temporarily that their mean is \(o(S)\).
\Cref{com:variation} then gives the same bound for center exceptions. We will
show that bit disagreement on $\mathcal C$ is also small at edge pairs at a point. Once both
variations for edge pairs at points and centers are small, expansion forces the edge bit itself to
be almost constant at nearly every level, contradicting \(S=\int f\).

\begin{proof}[Proof of \Cref{com:rare-positive}]
If the conclusion fails, restrict to a subsequence on which
\(\E_\mu H_p^*=o(S)\). By \eqref{com:center-rare}, also
\(\E_\kappa H_q^*=o(S)\).

For a probability measure \(\nu\) on states, let \(\nu_A,\nu_{B_1}\) denote
its coordinate marginals, with squared norms equal to the sums of squares of
their atom masses. Define
\begin{equation}\label{com:T}
 T=\int_{\mathcal C}\E_\mu\left[
       \norm{(\nu_{H(p)})_A}_2^2+
       \norm{(\nu_{H(p)})_{B_1}}_2^2-1\right].
\end{equation}
These norms count coordinate matches between two independent predicted
points with the same history. This correction is needed when the good neighbors of a point
have two states in one row or column and a third state outside it, while
all three sets \(R_y\) have nearly full prediction mass. Their
intersection then concentrates \(\nu_{H(p)}\) on a single state, making
the integrand in \eqref{com:T} close to one. The score for an edge pair at a point alone has only the general lower bound \(-1/2\) in this case;
adding the two gives a positive lower bound. We will compare \(\mathcal D_P+T\)
with \(C_P\) using \eqref{com:second-slice-score}. To retain the upper
bound under the present assumption on exception mass on $\mathcal G$, we first prove
\begin{equation}\label{com:prediction-score}
 T\le o(S).
\end{equation}

Choose a history with its marginal probability, then two independent points
\(p',p''\) from \(\mu\) conditional on that history. This is one fixed
coupling throughout all levels, and each point marginal is \(\mu\). At a
level, the expected value of
\[
 F(p',p'')=
   \1_{\{A_{p'}=A_{p''}\}}+
   \1_{\{B_{1,p'}=B_{1,p''}\}}-1
\]
is the integrand in \eqref{com:T}.

For a distinct pair put \(Z=z_{p'}-z_{p''}\) and
\(W=w_{p'}-w_{p''}\), both nonzero. At a place with level interval
\([l,u]\), the endpoint conditions from \eqref{ht:range} say
\(l\le-\log|Z|_v\) and \(u\ge\log|W|_v\).
Put \(a=-\log|Z|_v\), \(b=\log|W|_v\), and let \(\upsilon_v\) be
\(\log\sqrt2\) at infinity and zero otherwise.
An \(A\)-match requires \(s\le a+\upsilon_v\); a \(B_1\)-match requires
\(s\ge b-\upsilon_v\). Since \(a\ge l\) and \(b\le u\), the two
matching indicators have integrals at most \(\min(u,a+\upsilon_v)-l\)
and \(u-\max(l,b-\upsilon_v)\), respectively. Including the constant
term \(-1\) in \(F\) therefore gives
\begin{align*}
 \int_l^u F(p',p'')\dd s
 &\le [\min(u,a+\upsilon_v)-l]
       +[u-\max(l,b-\upsilon_v)]-(u-l)\\
 &=\min(u,a+\upsilon_v)-\max(l,b-\upsilon_v)\\
 &\le-\log|Z|_v-\log|W|_v+2\upsilon_v.
\end{align*}
In particular, if \(a+\upsilon_v<b-\upsilon_v\), then \(F=-1\)
throughout the gap \((a+\upsilon_v,b-\upsilon_v)\subseteq[l,u]\).
Thus the estimate also covers disjoint windows, when the bound is negative.
The product formula, and total archimedean weight one, show that the
integral over all levels is at most an absolute constant.

At a level in $\mathcal G$ where neither point is exceptional, their dominant
coordinate labels match, so \(F\ge0\). In general at such a level
\[
 F\ge-\1_{\{p'\text{ exceptional}\}}-
         \1_{\{p''\text{ exceptional}\}}.
\]
The contribution on levels in $\mathcal B$ is at least \(-W_b\). Removing all
levels outside $\mathcal C$ from the full integral consequently gives, for each
distinct pair,
\[
 \int_{\mathcal C}F(p',p'')
       \le O(1)+H_{p'}^*+H_{p''}^*+W_b.
\]

On the diagonal the integral on $\mathcal C$ is exactly \(W_c\). The
selected coupling is the pair law described after
\eqref{pred:point-collision}, so its diagonal probability satisfies
\[
 \Prob(p'=p'')\le |H(P)|\max_p\mu(p)=t^{-1+o(1)}.
\]
Here \(|H(P)|\le L^{2K}=t^{o(1)}\) and the marginal degree bounds were
proved in \Cref{sec:incidence}. Since
\(W_c\le k^{1/16}S\) and \(k=t^{o(1)}\), its contribution is \(o(S)\).
Average the bound for distinct pairs, using the two marginals \(\mu\),
\(W_b=o(S)\), and \(S\to\infty\). This proves
\(T\le2\E_\mu H_p^*+o(S)=o(S)\).

At a level in $\mathcal C$ declare an ordinary edge bad if its test is inconsistent
or if
\(\nu_{H(p)}(R_{\mathbf y_q})<1-k^{-2}\). Its integrated probability is \(o(S)\),
because Markov's inequality bounds the second contribution by
\(k^2\int_{\mathcal C}\E_E[1-\nu_{H(p)}(R_{\mathbf y_q})]=o(S)\).
Fix a point \(p\) and this level. Write \(b_p\) for its conditional bad-edge
probability, \(\nu=\nu_{H(p)}\), \(\zeta=k^{-2}\), and
\[
 \tau(\nu)=\norm{\nu_A}_2^2+\norm{\nu_{B_1}}_2^2-1.
\]
Let \(\mathcal Y_p\) be the states of its good neighbors. Every such state
has \(\nu(R_y)\ge1-\zeta\), and the actual point state belongs to each
\(R_y\), by consistency. Let \(d_P(p)\) be the score on $\mathcal C$ at this
edge pair at a point, using its full bit probability \(\rho_p\). We claim
\begin{equation}\label{com:second-slice-score}
 d_P(p)+\tau(\nu)
 \ge\frac\varpi4\,2\rho_p(1-\rho_p)-C(b_p+\zeta).
\end{equation}
As before, a pair with a bad slot costs only \(O(b_p)\); its probability is
at most \(2b_p\). The classification by intersections of the sets \(R_y\) proves the claim as
follows.

If \(\mathcal Y_p\) is empty, \(b_p=1\), while \(d_P(p)\ge-1/2\) and
\(\tau(\nu)\ge-1\), so the claim holds for an absolute \(C\).

If \(\mathcal Y_p=\set{(a,b)}\), put
\(r=\nu_A(a)\), \(s=\nu_{B_1}(b)\). Since
\(r+s\ge\nu(R_{(a,b)})\ge1-\zeta\),
\[
 \tau(\nu)\ge r^2+s^2-1
        \ge\tfrac12(1-\zeta)^2-1\ge-\tfrac12-\zeta.
\]
Two good neighbors have the same state and bit, so
\(d_P(p)\ge1-\varpi-O(b_p)\). Their sum has the positive lower bound
needed for \eqref{com:second-slice-score}.

If two states share a row, the intersection of their sets \(R_y\) is in that row, of
\(\nu\)-mass at least \(1-2\zeta\). Thus
\(\tau(\nu)\ge(1-2\zeta)^2-1\ge-4\zeta\).
If every state is in that row, the actual point lies there too, since it
belongs to the sets \(R_y\) for both distinct states. All good edges consequently
have bit \(1\). Their full-law variation is at most \(2b_p\), and
\(d_P(p)\ge-2b_p\), proving the claim in this case. Columns give bit \(0\)
and the same estimates. If there is a third state outside the row or column,
the three sets \(R_y\) have intersection of size at most one and mass at least
\(1-3\zeta\). It must be a singleton for large \(k\). Both marginals then
have an atom of at least that mass, giving
\[
 \tau(\nu)\ge2(1-3\zeta)^2-1\ge1-12\zeta.
\]
Together with \(d_P(p)\ge-1/2\), this is more than sufficient.

Finally suppose no coordinate is shared by distinct states in
\(\mathcal Y_p\). Three such states are impossible, since the actual point
belongs to all their sets \(R_y\) and their intersection is empty. In the
remaining case there are two states differing in both coordinates. The two
opposite corners carry \(\nu\)-mass at least \(1-2\zeta\). Each coordinate
marginal therefore has at least that mass on a set of size two, so
\[
 \tau(\nu)\ge
  2\cdot\frac{(1-2\zeta)^2}{2}-1\ge-4\zeta.
\]
The actual point is one of those corners. Its bit determines the good
neighbor state; hence \eqref{com:score-polynomial} gives
\(d_P(p)\ge\varpi/2-O(b_p)\). This completes the proof of
\eqref{com:second-slice-score}.

Average \eqref{com:second-slice-score} over \(p\) and integrate on levels in $\mathcal C$. The bad-edge integral is \(o(S)\), and
\(\zeta W_c\le k^{-31/16}S=o(S)\). The result is
\[
 (\text{contribution on $\mathcal C$ to }\mathcal D_P)+T
       \ge\frac\varpi4 C_P-o(S).
\]
The contribution outside $\mathcal C$ to \(\mathcal D_P\) is at least \(-o(S)\) by
\Cref{bud:rare-lower}, under the present assumptions on both exception
means. Therefore
\[
 \frac\varpi4 C_P-o(S)
 \le \mathcal D_P+T\le10\delta^2\varpi C_P+o(S),
\]
and again \(10\delta^2<1/4\) yields \(C_P=o(S)\).

To finish, let \(V_P,V_Q\) be the variations for edge pairs at points and centers at any fixed
level. Assign to each vertex its conditional majority bit, breaking ties
in a fixed way. The probability that the bit on an incident edge differs
from this label is at most the corresponding average variation, because
\(\min(\rho,1-\rho)\le2\rho(1-\rho)\). Thus the two endpoint labels disagree
on an edge with probability at most \(V_P+V_Q\). When that sum is small,
the matching-label consequence of expansion in \Cref{inc:matching} gives
one bit label with total marginal exception probability
\(O(V_P+V_Q)\). The edge bit differs from it with probability
\(O(V_P+V_Q)\) as well. When the sum is not small the same bound follows
from \(f\le1/2\), with a larger absolute constant. Consequently
\[
 f\le C(V_P+V_Q)
\]
uniformly for all sufficiently large \(t\). Integrating on levels in $\mathcal C$
gives \(\int_{\mathcal C}f\le C(C_P+C_Q)=o(S)\).

On levels in $\mathcal G$, the minority-bit probability is bounded by the two
marginal exception probabilities and the inconsistency probability, as in
\Cref{bud:rare-moments}. Its integral is \(o(S)\) under the present
assumptions. Levels in $\mathcal B$ contribute at most \(W_b/2=o(S)\). Thus
\(S=\int f=o(S)\), a contradiction. This proves
\eqref{com:rare-positive-eq}.
\end{proof}

The positive exception mass on $\mathcal G$ is a first-moment statement. To extract many
pairs at a single height scale, \Cref{sec:extraction} next controls its
tail of the exception mass and removes errors in measure weighted by exception mass before
selecting a band. Throughout that selection, the bound $C_Q=o(S)$ for the bit-disagreement integral in experiment $Q$ remains available under the original law of an edge pair at a center.
